\section{実験計画法4；一要因モデル(Within)}
\label{lesson14}

\subsection{授業内容}
\subsubsection{科目の中でのこのコマの位置づけ}
\begin{tabular}[t]{ccccccc}
    記述統計[4] & 確率[2] & モデル[8] & 推測・推定[1] & 意思決定[0] & 実践[2]
\end{tabular}
\subsubsection{概要}
再び一要因に戻るが，今度は群内(Within)計画の場合で考える。群内計画の場合は，事前-事後や第I期，II期，III期など反復測定をしている場合になり，個体の識別ができていることになる。そのことで個人差を考慮することができるようになり，より精緻に誤差を取り除くことができるとも言える。効果の比較対象とするべき誤差の計算方法，数式表現などBetweenデザインと異なるところに注意しながら理解を深める。
\subsubsection{コマ主題細目(箇条書き)}
\begin{description}
\item[Within計画のデータ] 一要因Withinの実験計画法モデルを考えるために，数値例を導入する。ここでの数値例は\citet{Yamada2004}のP.178の表7.5.1である。可視化に当たっては，個人が識別できていることから散布図ではなく折れ線グラフで描画するのが適切である。可視化されたグラフを見ながら，どのような数式表現が可能かを考える。

\item[モデルと誤差] 個人差を個人ごとの平均と考え，数式で表現する。次に各データから個人平均を取り除いて数値で考えてみる。こうすることで，誤差がどこに現れているかを直観的に考えることができる。これらのプロセスを通じて数式を完成させ，また個々のセルにあらわれる誤差を算出する。
\begin{flushright}
    $\to$ 計算の手順は\citet{Yamada2004}のP.178--183.
\end{flushright}

\item[BetweenとWithinの違い] 同じようなデータであっても，個人が識別されていることで個人差を算出することができた。これはBetweenデザインでは個人差も要因計画では考慮できない誤差として捉えていたことを意味する。逆に，WithinデザインはBetweenデザインよりも誤差を精緻に見積もることができるとも言える。心理学における実験計画は一般にWithinデザインにする方が有利であるとも言えるが，Withinデザインは同じ個体に反復測定を行うことでもあり，実験上運用上のコストは大きいことに留意する。
\begin{flushright}
    $\to$ 計算の手順は\citet{Yamada2004}のP.179，図7.5.1の平方和の分解図式
\end{flushright}

\item[誤差の分布，個人差の分布] ガウスの考えた誤差論は，誤差は真値の周りで正規分布することを表していた。実験計画法における誤差も，理論値の周りで正規分布すると考えられる。一方で，さまざまな複合的な原因から影響されるものは自然と正規分布するため，個々人の生理・生態学的特徴も正規分布に従う。Withinデザインは複数の正規分布が含まれる混合モデルである。

\end{description}
\subsubsection{キーワード}
\begin{itemize}
\item Withinデザイン
\item 個人差
\item 階層線形モデル
\item 変量効果と固定効果
\end{itemize}
\subsection{授業情報}
\paragraph{コマの展開方法}
講義
\paragraph{標準シラバスにおける位置づけ}
科目番号4；心理学研究法；2.データを用いた実証的な思考方法；C データの統計的記述\\
科目番号5；心理学統計法；1.心理学で用いられる統計手法；G 推測統計:代表値と散布度をめぐって(1)
\subsubsection{予習・復習課題}
\paragraph{予習}
実験計画の用語を復習し，どういった実験計画が群内Withinデザインになるのかを考える。心理学系論文の中から，群内デザインの実験を行なっている例を探して読んでみるのも予習になる。
\paragraph{復習}
Withinデザインにおいて個人差を算出するための手続きを，もう一度自分で確認しておくと良い。計算にあたっては身近な数値例を探すと良いが，個人が特定できるようなデータは見つけにくいかもしれない。そのような場合は，Betweenデザインの数値例を，借りにWithinデザインで得られた数字であるとみなして，計算上の練習を行うのも効果的である。
また今回，最後に後期の確率モデルにつながる「混合モデル」の考え方が示された。混合モデルについては，階層線形モデル，マルチレベルモデルなど複数の呼称が存在するので，参考資料を探す場合はそれぞれのキーワードで探すと良いだろう。